\documentclass{situationist-preprint}

\title{Utopia and Reasonable Pluralism}
\msid{POL-21-a}
\msfamily{POL-21}
\msdiscipline{Political Philosophy}
\mstype{Principal manuscript}
\mskeywords{utopia; reasonable pluralism; value pluralism; Berlin; Nozick; Rawls}

\begin{document}
\maketitle

\begin{abstract}
Is a perfect society possible? We argue that it is not: every utopia conflicts with some reasonable conception of the good. A society is perfect only if no reasonable conception of the good gives those who hold it reason to judge some feasible alternative better. The good is plural, its components cannot be fully realized together, and every society settles the trade-offs among them through its institutions and its culture. Reasonable conceptions of the good weigh the components differently, and for every settlement some reasonable conception ranks another higher. Hence no society is perfect. We consider Nozick's framework for utopia, which seeks to escape the argument by letting each community realize its own conception, and argue that the framework is itself one settlement among others. The result is consistent with the legitimacy of a liberal state, which requires justifiability to all reasonable persons, not perfection by every reasonable standard.
\end{abstract}

\section{The problem}

\citet{More1516} gave the genre its name and its ambiguity: the good place that is no place. The utopian tradition, from Plato's \emph{Republic} to the nineteenth-century socialists, held that a society could be arranged so that nothing of importance was lost. Its critics held that the attempt to realize a perfect society licenses the sacrifice of actual persons to an ideal. \citet{Popper1945} argued against utopian social engineering on epistemic grounds; \citet{Berlin1958,Berlin1990} argued against it on grounds of value: since human goods are many and some are incompatible, the idea of a final harmony in which all are realized is incoherent. \citet[Part~III]{Nozick1974} replied that a framework of voluntary communities, in which each may pursue its own vision, is the one utopia compatible with diversity. \citet{Rawls1993} observed that a free society will always contain a plurality of reasonable comprehensive doctrines, and \citet{Rawls1999} defended a \enquote{realistic utopia} that is explicitly not a perfect society. We make Berlin's argument precise and show that Nozick's reply does not escape it.

\section{Framework}

\begin{definition}[Perfect society]\label{d:perfect}
A feasible society is \emph{perfect} if no reasonable conception of the good gives those who hold it reason to judge some other feasible society better.
\end{definition}

\begin{definition}[Settlement]\label{d:settlement}
A society's \emph{settlement} is the way its institutions and culture weigh the components of the good against one another where they cannot all be realized together: how much they favor stable community over mobility, achievement over leisure, autonomy over shared tradition.
\end{definition}

\begin{assumption}[Plurality and conflict]\label{a:conflict}
The good has several irreducible components, and no feasible society realizes all of them fully; every feasible society has a settlement.
\end{assumption}

\begin{assumption}[Reasonable divergence]\label{a:diverge}
For every settlement, there is a reasonable conception of the good that weighs the components differently, and that ranks some feasible society with a different settlement above it.
\end{assumption}

Assumption~\ref{a:conflict} is the conclusion of ETH-01, on which the good has seven irreducible components, together with the observation that some of them compete for scarce time and resources and that some are realized only in social forms that exclude one another. A tightly knit community sharing a way of life realizes goods of relationship that a mobile society of chosen associations does not, and the reverse holds for goods of autonomy; abundance does not remove such conflicts, since they concern social form rather than resources. Assumption~\ref{a:diverge} expresses reasonable pluralism in Rawls's sense: because of the burdens of judgment, conscientious and informed persons reach different weightings of the goods \citep[Lecture~II]{Rawls1993}, and these weightings differ, among other things, in the settlements they favor.

\section{The argument}

\begin{theorem}\label{t:main}
No feasible society is perfect.
\end{theorem}

\begin{proof}
Let $S$ be a feasible society. By Assumption~\ref{a:conflict}, $S$ has a settlement. By Assumption~\ref{a:diverge}, some reasonable conception of the good ranks a feasible society with a different settlement above $S$, and so gives those who hold it reason to judge that society better. By Definition~\ref{d:perfect}, $S$ is not perfect.
\end{proof}

The theorem says more than that no actual society is perfect. It says that every utopia, every description of a society proposed as perfect, conflicts with some reasonable conception of the good, because a utopia that is described in enough detail to be a society has a settlement.

\section{Objections}

\begin{objection}[The framework for utopia]
Nozick's framework does not impose a settlement: it lets each community choose its own, and lets individuals move among them. Every reasonable conception can be realized in some community, so none has reason to prefer an alternative.
\end{objection}

\begin{reply}
The framework is itself a settlement. It gives priority to exit and voluntary association, and so to the components of the good that choice serves, over those that require that a way of life be shared by a whole society and secure from defection. Conceptions on which a common culture, a shared faith or a stable polity is a central good are not irrational or unreasonable, and they rank above the framework a society in which that good is realized for all. Nozick's framework may be the best settlement; it is not a society without loss.
\end{reply}


\begin{objection}[Acceptance]
A reasonable person would accept a fair compromise among conceptions of the good. A society built on such a compromise is perfect, since no reasonable person rejects it.
\end{objection}

\begin{reply}
To accept a compromise is not to judge it best. Definition~\ref{d:perfect} asks whether a reasonable conception gives reason to judge another society better, and a person may have such reason while accepting that the compromise is fair. This is the distinction between legitimacy and perfection. On POL-02, a state is legitimate if it can be justified to all reasonable persons subject to it; a legitimate state is one that reasonable persons can accept, not one that each judges best. Theorem~\ref{t:main} is therefore compatible with the existence of legitimate states, and indeed with a realistic utopia in Rawls's sense, which aims at what is just and feasible rather than at a society without loss.
\end{reply}

\section{Conclusion}

Because the good is plural, every society settles the conflicts among its components; because reasonable conceptions of the good diverge, every settlement is ranked below some alternative by a reasonable conception. No society is perfect, and every utopia conflicts with some reasonable conception of the good. What remains possible is a society that all reasonable persons can accept.

\begin{thebibliography}{99}\small
\bibitem[Berlin(1958)]{Berlin1958} Berlin, I. (1958). \emph{Two Concepts of Liberty}. Clarendon Press.
\bibitem[Berlin(1990)]{Berlin1990} Berlin, I. (1990). \emph{The Crooked Timber of Humanity: Chapters in the History of Ideas}. Ed. H. Hardy. John Murray.
\bibitem[More(1516)]{More1516} More, T. (1516). \emph{Utopia}.
\bibitem[Nozick(1974)]{Nozick1974} Nozick, R. (1974). \emph{Anarchy, State, and Utopia}. Basic Books.
\bibitem[Popper(1945)]{Popper1945} Popper, K.~R. (1945). \emph{The Open Society and Its Enemies}. 2 vols. Routledge.
\bibitem[Rawls(1993)]{Rawls1993} Rawls, J. (1993). \emph{Political Liberalism}. Columbia University Press.
\bibitem[Rawls(1999)]{Rawls1999} Rawls, J. (1999). \emph{The Law of Peoples}. Harvard University Press.
\end{thebibliography}

\end{document}
